\begin{sol}[Barra pivotada con partícula deslizante y resorte]

\subsubsection*{Pregunta 1: Momento de inercia de la barra}

Con densidad lineal $\rho=M/L$, el momento de inercia respecto al centro de masa es
$$\bar{I}_z = \int_{-L/2}^{L/2} s^2\,\frac{M}{L}\,ds = \frac{M}{L}\cdot\frac{2}{3}\left(\frac{L}{2}\right)^3 = \frac{1}{12}ML^2.$$
Por el teorema de Huygens-Steiner, trasladando al extremo $O$ (a distancia $L/2$):
$$I_z^O = \bar{I}_z + M\left(\frac{L}{2}\right)^2 = \frac{1}{12}ML^2 + \frac{1}{4}ML^2 = \boxed{\tfrac13 ML^2}.$$

\subsubsection*{Pregunta 2: Ecuaciones de movimiento}

\textbf{Cinemática.} Se usa una base polar $\{\vec{e}_r,\vec{e}_\theta\}$ solidaria a la barra
($\vec{e}_r$ a lo largo de la barra). La partícula está en $\vec{r}_m=r\,\vec{e}_r$, de modo que
$$\vec{a}_m = (\ddot{r}-r\dot\theta^2)\,\vec{e}_r + (2\dot{r}\dot\theta + r\ddot\theta)\,\vec{e}_\theta.$$
Con $\theta$ medido desde la vertical descendente, la gravedad se proyecta como
$g\cos\theta\,\vec{e}_r - g\sin\theta\,\vec{e}_\theta$.

\textbf{Newton sobre la partícula.} La barra lisa ejerce sólo una normal $\vec{N}=-N\vec{e}_\theta$;
el resorte da $-K(r-r_0)\vec{e}_r$:
$$\vec{e}_r:\quad -K(r-r_0) + mg\cos\theta = m(\ddot{r}-r\dot\theta^2),$$
$$\vec{e}_\theta:\quad -N - mg\sin\theta = m(2\dot{r}\dot\theta + r\ddot\theta).$$
De la primera ecuación se despeja directamente
$$\boxed{\;\ddot{r} = -\frac{K}{m}(r-r_0) + g\cos\theta + r\dot\theta^2\;}$$

\textbf{Euler sobre la barra respecto a $O$.} La gravedad de la barra (en $L/2$) y la reacción
$N$ de la partícula (a distancia $r$) generan el torque:
$$I_z^O\,\ddot\theta = -\frac{L}{2}Mg\sin\theta + rN,\qquad I_z^O=\tfrac13 ML^2.$$
Sustituyendo $N = -mg\sin\theta - m(2\dot{r}\dot\theta + r\ddot\theta)$ y agrupando:
$$\boxed{\;\ddot\theta = \frac{-g\sin\theta\left(\tfrac{L}{2}M + rm\right) - 2mr\dot{r}\dot\theta}{\tfrac13 ML^2 + mr^2}\;}$$

\end{sol}
